\section{Related work}
\label{sec:related-work}

The test sits where four bodies of work meet: the supervisory documents that define the measure and state the practice, the models of the deposit whose parameters the test reads, the empirical studies of how reported risk measures and accounting assumptions behave near a threshold, and the statistical methods for detecting a change and for testing the association of two events. Each is characterised below by what it assumes, what it provides, and where it stops short of the question of this paper.

\subsection{Supervisory standards and the review of the failure}
\label{sec:rw-supervisory}

The Basel standard on interest rate risk in the banking book of April 2016 sets the standardised framework in which the loss of economic value under six prescribed shocks is measured, with caps on the share of a non-maturity deposit that can be treated as stable core and on the average maturity of that core, by type of depositor, and an outlier test that compares the loss with a fraction of Tier~1 capital \citep{bcbs2016}. The text now lives in chapter SRP31 of the Basel Framework \citep{basel2019srp31}, and the sizes of the shocks were recalibrated in July 2024 \citep{bcbs2024}. These documents fix the measure and the limits of the behavioural assumptions, and they assume that the bank, or its validator, controls what the assumptions are at each date. They provide no check on the history of the assumptions. The Federal Reserve's review of the failure of Silicon Valley Bank is the source of the practice that motivates this work: it states that the bank breached its long-term interest rate risk limits on and off since 2017 and that in April 2022 it made counterintuitive assumptions about the duration of deposits to address a breach \citep{fed2023}. A finding of that kind is made once, by a supervisor with access to the model files and to the minutes. It assumes that someone has read the history of the assumption next to the history of the limit, and it does not say how to do that when the history is a table of fitted parameters and a column of ratios. The test of this paper is a procedure for that reading.

\subsection{Models of non-maturing deposits}
\label{sec:rw-deposits}

The parameters that the test reads come from models of deposits. Hutchison and Pennacchi measure the rents and the interest rate risk of retail bank deposits in imperfect financial markets \citep{hutchison1996}. Jarrow and van Deventer value demand deposits and credit card loans by arbitrage-free arguments and derive hedges \citep{jarrow1998}, and Kalkbrener and Willing build a model of non-maturing liabilities for risk management \citep{kalkbrener2004}. On the empirical side, Drechsler, Savov and Schnabl show that when the policy rate rises banks widen the spreads they charge on deposits and deposits leave the banking system, more so where deposit markets are concentrated \citep{drechsler2017}, and that the resulting deposit franchise lets banks transform maturity without bearing the interest rate risk that the assets would imply \citep{drechsler2021}. Jiang and co-authors study the mark-to-market losses of US banks that followed the monetary tightening and their relation to runs by uninsured depositors \citep{jiang2024}. These works assume that deposit behaviour is a stable function of the market rate or of the competitive structure, and they provide the quantities, a pass-through, a core share, a duration, that a bank must assume. None of them concerns the timing of a change in an assumed quantity relative to a limit, and the papers on the franchise imply that the assumed duration is hard to pin down from data, which is the condition under which a change can be defended after the fact.

\subsection{Manipulation of assumptions and of regulatory measures}
\label{sec:rw-manipulation}

A literature in accounting and banking asks whether reported quantities and discretionary assumptions behave differently where an incentive exists. Burgstahler and Dichev document that the distribution of reported earnings has fewer small losses and more small profits than a smooth distribution would give, and attribute it to the management of earnings to avoid decreases and losses \citep{burgstahler1997}. Healy and Wahlen review the evidence on which accruals are used to manage earnings, how large and how frequent the management is, and whether it affects the allocation of resources \citep{healy1999}. Bergstresser, Desai and Rauh study the assumptions that firms make about their pension plans and relate them to earnings manipulation and managerial incentives \citep{bergstresser2006}, which is a published analogue of an assumption chosen to fit a target. In banking, Ahmed, Takeda and Thomas re-examine whether loan loss provisions are used to manage capital and earnings \citep{ahmed1999}; Mariathasan and Merrouche study the manipulation of the risk weights that banks using internal models report \citep{mariathasan2014}; Begley, Purnanandam and Zheng study the strategic underreporting of bank risk \citep{begley2017}; and Plosser and Santos relate the inconsistency of the risk estimates of different banks for the same exposure to the banks' incentives \citep{plosser2018}. McCrary's density test checks for manipulation of a variable around a threshold by looking for a discontinuity in its density \citep{mccrary2008}, and Strathern's account of ratings in the audit of universities is the source of the general warning that a measure that becomes a target ceases to measure \citep{strathern1997}. These studies assume a cross-section of many firms or banks, or a long panel, so that the behaviour near the threshold can be estimated against the behaviour away from it. They provide evidence at the level of a population. They cannot be applied to one bank with one limit and a dozen quarters, which is the situation of a validator, and they do not test the timing of one assumption against the dates of one breach.

\subsection{Change detection and tests of association}
\label{sec:rw-statistics}

Sequential and retrospective methods detect that a parameter changed. The cumulative sum scheme of Page signals a shift in the mean of a sequence \citep{page1954}, the cumulative sum of recursive residuals tests the constancy of a regression relationship over time \citep{brown1975}, Chow's test compares the coefficients of a regression before and after a date known in advance \citep{chow1960}, and the supremum tests of Andrews and the procedures of Bai and Perron locate one or several breaks at dates that are not known \citep{andrews1993,bai1998}. They assume a stochastic model for the noise around the parameter, and their power comes from many observations on each side of the break. They are unconditional on an external flag, with the exception of Chow's test, which takes the date as given and which, for a constant-mean model with the breach indicator as the only regressor, is the two-sample comparison of the mean change at breach dates with the mean change elsewhere; the unequal-variance form of that comparison is Welch's \citep{welch1947}. Fisher's exact test of the association of two binary events in a two-by-two table \citep{fisher1935} and the rank test of Mann and Whitney for the size of the values in two groups \citep{mann1947} are the other tests that can be applied to the question, here the association of the event ``the parameter changed'' with the event ``a limit was breached'' and the comparison of the size of the changes at breach dates with those elsewhere. They assume that the dates are exchangeable under the null. They provide a p-value, and at the lengths of history of a validation they cannot reject unless the evidence is strong, as Section~\ref{sec:theory} shows.

\subsection{Where the proposed test stands}
\label{sec:rw-position}

The proposed test is not a change-point method, since it does not model the noise of the parameter, and it is not a population test of manipulation, since it uses one bank. It is a deterministic rule of coincidence on a within-bank history, closer to an exception report than to an hypothesis test, with the property that its decision on a stated parameter can be read without a model. The published tests above are its statistical analogues, and the engine's own test of the stability of the replicating-portfolio weights, which follows the technique that the standard describes \citep{bcbs2016}, is the existing check of the behavioural parameters; it ignores the breach flags. Section~\ref{sec:experimental-design} runs the Fisher, Mann--Whitney and Welch tests, the stability test and a breach-agnostic count on the same histories as the proposed test, so that the comparison is of the procedures and not of what their authors reported. What the literature does not provide, and this paper adds, is the formal statement of the coincidence test, its false-flag and detection probabilities in closed form, its behaviour on the two kinds of parameter a bank holds, and a tolerance rule that restores discrimination when the parameter is re-estimated at every date.
